\documentclass[11pt,a4paper]{book}
\usepackage[margin=1in]{geometry}
\usepackage[utf8]{inputenc}
\usepackage{setspace}
\usepackage{parskip}
\usepackage{csquotes}
\usepackage{microtype}
\usepackage{hyperref}
\usepackage{fancyhdr}
\usepackage{xcolor}
\usepackage{hyperref}
\usepackage{fontspec}
\usepackage{etoolbox} % for patchcmd
\usepackage[compact]{titlesec}
\usepackage{minted}
\usepackage{enumitem}
\usepackage{amsmath}
\usepackage{listings}
\usepackage{xcolor}
\usepackage{multicol}
\usepackage{graphicx}
\usepackage{tikz}
\usepackage{xcolor}
\usepackage{etoolbox}
\usepackage{framed}
\usepackage{float}
\usepackage{pdfpages}
\usepackage{metalogo}
\usepackage{setspace}

\setstretch{1.15}

\definecolor{accentgray}{RGB}{80, 80, 80}

\onehalfspacing

\pagestyle{fancy}
\fancyhf{}
\renewcommand{\headrulewidth}{0pt}
\fancyfoot[C]{\thepage} % Put page number in center footer

\hypersetup{
    colorlinks=true,
    linkcolor=black,
    urlcolor=accentgray,
}

\newfontfamily{\LatinModernRoman}{Latin Modern Roman}

\setmainfont[
    Path=fonts/,
    UprightFont = *-Regular,
    BoldFont    = *-Bold,
]{Arial}

\setmonofont[
    Path=fonts/,
    Scale=0.9,
    UprightFont = *-Regular,
    BoldFont    = *-Bold,
]{ComicMono}

\hyphenpenalty=9999
\exhyphenpenalty=10000
\sloppy
\renewcommand{\arraystretch}{1.5}
\setlength{\emergencystretch}{4em}
\setlength{\parskip}{1em}
\AtBeginEnvironment{tabular}{\fontsize{10}{12}\selectfont}
\titleformat{\chapter}[display]{\normalfont\huge\bfseries}{\chaptertitlename\ \thechapter}{-0.5em}{\Huge}
\titlespacing*{\chapter}
  {0pt}    % left margin
  {-5em}    % space before
  {-0.5em}    % space after
\setlength{\parindent}{0pt}

\renewcommand{\contentsname}{Table of Contents}

\lstdefinestyle{text}{
    basicstyle=\ttfamily\color{black},
    breaklines=true,
}

\definecolor{white}{HTML}{ffffff}
\definecolor{purple}{HTML}{a68cd9}

\makeatletter

% PATCH: CHAPTER
\patchcmd{\chapter}{\if@openright\cleardoublepage\else\clearpage\fi}{}{}{}
\let\old@chapter\chapter
\renewcommand{\chapter}{%
  \ifdim\pagetotal=0pt % remove margin if it's at top of page
    \vspace*{-2.5em}
  \fi
  \@ifstar
    {\old@chapter*}
    {\old@chapter}%
}
\titleformat{\chapter} % or [block]
  {\normalfont\LARGE\bfseries}
  {\thechapter}{0.625em}{}
\titlespacing*{\chapter}
  {0pt}    % left margin
  {0em}    % space before
  {-0.5em}    % space after

\makeatother

\usemintedstyle{tango}
\newcommand{\inputmintedstyled}[2]{%
  \inputminted[
    fontsize=\small,
    fontfamily=tt,
    bgcolor=white,
    rulecolor=\color{purple},
    linenos,
    numbersep=10pt,
    frame=single,
    framesep=2mm,
    breaklines,
    tabsize=4,
    breakafter=-/\{,
    breakbefore=\,\\
  ]{#1}{#2}%
}

\newcommand{\inputmintedstyledtwocolumns}[2]{%
  \begin{multicols*}{2}
  \inputminted[
    fontsize=\small,
    fontfamily=tt,
    bgcolor=white,
    rulecolor=\color{purple},
    linenos,
    numbersep=0pt,
    frame=single,
    framesep=2mm,
    breaklines,
    tabsize=4,
    breakafter=-/\{,
    breakbefore=\,\\
  ]{#1}{#2}
  \end{multicols*}%
}

\newcommand{\quotefont}{\rmfamily}

\newcommand*\quotesize{80} % if quote size changes, need a way to make shifts relative
% Make commands for the quotes
\newcommand*{\openquote}
   {\tikz[remember picture,overlay,xshift=-2.5ex,yshift=-4.5ex]
   \node (OQ) {\quotefont\fontsize{\quotesize}{\quotesize}\selectfont``};\kern0pt}

\newcommand*{\closequote}[1]
  {\tikz[remember picture,overlay,xshift=2.5ex,yshift={#1}]
   \node (CQ) {\quotefont\fontsize{\quotesize}{\quotesize}\selectfont''};}

% select a colour for the shading
\colorlet{shadecolor}{gray!10}

\newcommand*\shadedauthorformat{\emph} % define format for the author argument

% Now a command to allow left, right and centre alignment of the author
\newcommand*\authoralign[1]{%
  \if#1l
    \def\authorfill{}\def\quotefill{\hfill}
  \else
    \if#1r
      \def\authorfill{\hfill}\def\quotefill{}
    \else
      \if#1c
        \gdef\authorfill{\hfill}\def\quotefill{\hfill}
      \else\typeout{Invalid option}
      \fi
    \fi
  \fi}
% wrap everything in its own environment which takes one argument (author) and one optional argument
% specifying the alignment [l, r or c]
%
\newenvironment{shadequote}[2][l]%
{
\vspace{1em}
\authoralign{#1}
 \ifblank{#2}
   {\def\shadequoteauthor{}\def\yshift{-2ex}\def\quotefill{\hfill}}
   {\def\shadequoteauthor{\par\authorfill\shadedauthorformat{#2}}\def\yshift{-2ex}}
 \begin{snugshade}\vspace{0.5em}\begin{quote}\openquote}
{\shadequoteauthor\quotefill\closequote{\yshift}\end{quote}\vspace{0.5em}\end{snugshade}%
}

\begin{document}

\begin{center}
{\large\bfseries ITEC 50: Web Systems and Technologies}

{\large\bfseries Final Project --- Group Progress Update Report}

\end{center}

\vspace{1em}
\hrule
\vspace{1em}

\noindent\textbf{Project Title:} DeciMark

\textbf{Members:} Lyra Phasma
\vspace{1.5em}
\hrule
\vspace{1em}

\chapter{Project Overview}

\textbf{DeciMark} is a \textit{Bookmark Link Manager} --- a web application that enables users to save, organize, retrieve, and annotate URLs using the \textbf{Johnny.Decimal} classification system. The application is server-backed with a full database layer, and requires an account to use. No client-side installation is required --- users access DeciMark entirely through a browser, with all complexity contained on the server side.

Target users are students and professionals who need structured, consistent link organization without the entropy of free-form tagging systems. In later development stages, DeciMark is planned to be distributable as a portable installable --- \texttt{.exe}, \texttt{.msi}, or \texttt{.AppImage} --- for users who prefer self-hosting without a full server setup.

\chapter{Current Progress}

Conceptualization and competitive analysis are complete. Core design philosophy and feature differentiation from existing tools --- including Pinboard, linkding, and Raindrop.io --- have been fully documented. The Johnny.Decimal classification scheme has been selected and justified as the primary organizational model. The project scope has evolved from an initial static concept into a full server-backed implementation, with user accounts and a persistent database layer.

\chapter{Planned Tasks}

\begin{itemize}
    \item Implementation of the card-based UI with list-view toggle
    \item Johnny.Decimal ID field integration and range-based filtering
    \item Tag support and multi-attribute search functionality
    \item User account system and authentication
    \item Cloud synchronization of bookmark data across sessions
    \item Persistent database integration
    \item Final polish, accessibility improvements, and export functionality
    \item Long-term: portable distribution as \texttt{.exe}, \texttt{.msi}, and \texttt{.AppImage}
\end{itemize}

\chapter{Tools and Technologies}

\begin{itemize}
    \item \textbf{Frontend:} HTML, CSS, JavaScript
    \item \textbf{Backend:} Python, FastAPI
    \item \textbf{Database:} PostgreSQL, SQLModel, Alembic (migrations)
    \item \textbf{Authentication:} JWT, itsdangerous, bcrypt
    \item \textbf{Design reference:} Johnny.Decimal organizational system (\url{https://johnnydecimal.com})
\end{itemize}

\chapter{Challenges / Concerns}

Working as a sole member required managing all responsibilities --- research, design, documentation, and implementation --- independently. This was not without its silver linings, however. A concurrent hackathon project provided substantial hands-on experience in full-stack development, database architecture, backend frameworks, and deployment pipelines, all of which have meaningfully deepened my understanding of web systems at a level that directly informs the execution of DeciMark.

\chapter{Additional Notes}

It was not all in vain.

The hackathon experience --- undertaken in parallel with this course --- resulted in the development of \textbf{Pebble Recall}, a full-stack psychiatric care companion application built with Python, FastAPI, SQLModel, PostgreSQL, and Jinja2, stored at \url{https://github.com/phasmolysis-team/phasmolysis-pebble-recall} and deployed live at \url{https://pebble-recall.lyra-on.top}.

Pebble Recall placed \textbf{\#1 in judge's choice} and \textbf{\#1 in community vote} out of more than 140 competing teams, winning the \textbf{₱20,000 grand prize}.

The skills developed over the course of that project --- including database migration with Alembic, asynchronous backend architecture, authentication systems, and full deployment pipelines --- are not merely adjacent to DeciMark. They are directly applicable to it. The detour, it turns out, was not a detour at all.

DeciMark is currently in active development, with the backend implementation underway. The project repository is publicly available at \url{https://github.com/whinee/ITEC50-Finals}. A live deployment will follow as the frontend and backend are integrated.

To promote genuine learning while protecting both intellectual property and academic integrity, the DeciMark repository is multi-licensed:

\begin{itemize}
    \item \textbf{Mozilla Public License 2.0} --- applied to files derived from Pebble Recall, in compliance with that project's licensing terms
    \item \textbf{Lyra Phasma Academic Source License v1.0 (LPASL v1.0)} --- applied to all original source code. This license permits students to study and learn from the codebase, while explicitly prohibiting direct copying or unattributed submission as one's own academic work
    \item \textbf{MIT License} --- applied to all remaining files, including tooling and configuration
\end{itemize}

The LPASL v1.0 is an original license authored by the developer, designed specifically for the academic context: it encourages learning from open source without enabling academic dishonesty, and includes a sunset clause converting all restricted works to MIT terms on December 31, 2030.

\end{document}